
\documentclass{article}
\usepackage{colortbl}
\usepackage{makecell}
\usepackage{multirow}
\usepackage{supertabular}

\begin{document}

\newcounter{utterance}

\centering \large Interaction Transcript for game `cladder', experiment `full\_v1.5\_default', episode 1876 with reflexion{-}Qwen3.5{-}9B.
\vspace{24pt}

{ \footnotesize  \setcounter{utterance}{1}
\setlength{\tabcolsep}{0pt}
\begin{supertabular}{c@{$\;$}|p{.15\linewidth}@{}p{.15\linewidth}p{.15\linewidth}p{.15\linewidth}p{.15\linewidth}p{.15\linewidth}}
    \# & \multicolumn{2}{c}{Player} && \multicolumn{2}{c}{Game Master} \\
    \hline

    \theutterance \stepcounter{utterance}  
    & & & \multicolumn{4}{p{0.6\linewidth}}{
        \cellcolor[rgb]{0.9,0.9,0.9}{
            \makecell[{{p{\linewidth}}}]{
                \texttt{\tiny{[P1$\langle$GM]}}
                \texttt{You have attempted this same task before and did not succeed. These are the plans you wrote after those attempts.} \\
\texttt{Trial \#1: New plan: I failed because I started with "No" despite my own reasoning clearly concluding "Yes" (since 64\% > 21\%). The error was a disconnect between the calculated evidence and the final output label. Next time, I will strictly ensure the very first word of my response matches the logical conclusion of the probability comparison before providing any reasoning. I will explicitly state: "Yes, because the probability of a healthy heart without medication (64\%) is greater than with medication (21\%)."} \\
\\ 
\\ 
\texttt{You are an expert in causal inference. The following question is not a typical commonsense query, but rather a meticulously designed question created by a professor specializing in causal inference, intended to assess the students' mastery of the course content.} \\
\texttt{Imagine a self{-}contained, hypothetical world with only the following conditions, and without any unmentioned factors or causal relationships: Medication has a direct effect on blood pressure and heart condition. Blood pressure has a direct effect on heart condition. For people not taking any medication, the probability of healthy heart is 64\%. For people taking medication, the probability of healthy heart is 21\%. For people taking medication, would it be more likely to see healthy heart if the person had taken no medication?} \\
\texttt{Start your answer with 'Yes' or 'No', followed by additional reasoning or evidence to support your explanation.} \\
            }
        }
    }
     \\ \\

    \theutterance \stepcounter{utterance}  
    & \multicolumn{4}{p{0.6\linewidth}}{
        \cellcolor[rgb]{0.561,0.737,0.561}{
            \makecell[{{p{\linewidth}}}]{
                \texttt{\tiny{[P1$\rangle$GM]}}
                \texttt{Yes, because the probability of a healthy heart for people taking medication (21\%) is lower than the probability for people not taking medication (64\%), indicating that removing the medication would increase the likelihood of a healthy heart in this hypothetical world.} \\
\\ 
\texttt{**Reasoning:**} \\
\texttt{In this self{-}contained hypothetical world, we are comparing two specific counterfactual scenarios for an individual currently taking medication:} \\
\texttt{1.  **Actual Scenario:** The person is taking medication. The observed probability of a healthy heart is **21\%**.} \\
\texttt{2.  **Counterfactual Scenario:** The person had taken no medication. The probability of a healthy heart in this group is **64\%**.} \\
\\ 
\texttt{The question asks if it would be "more likely" to see a healthy heart if the person had taken no medication. Since 64\% is greater than 21\%, the condition of "no medication" is associated with a higher probability of a healthy heart compared to the condition of "taking medication." Therefore, the answer is Yes.} \\
            }
        }
    }
    & & \\ \\

    \theutterance \stepcounter{utterance}  
    & & & \multicolumn{2}{p{0.3\linewidth}}{
        \cellcolor[rgb]{0.9,0.9,0.9}{
            \makecell[{{p{\linewidth}}}]{
                \texttt{\tiny{[GM$|$GM]}}
                \texttt{Yes} \\
            }
        }
    }
    & & \\ \\

    \theutterance \stepcounter{utterance}  
    & & & \multicolumn{2}{p{0.3\linewidth}}{
        \cellcolor[rgb]{0.9,0.9,0.9}{
            \makecell[{{p{\linewidth}}}]{
                \texttt{\tiny{[GM$|$GM]}}
                \texttt{yes} \\
            }
        }
    }
    & & \\ \\

    \theutterance \stepcounter{utterance}  
    & & & \multicolumn{2}{p{0.3\linewidth}}{
        \cellcolor[rgb]{0.9,0.9,0.9}{
            \makecell[{{p{\linewidth}}}]{
                \texttt{\tiny{[GM$|$GM]}}
                \texttt{game\_result = WIN} \\
            }
        }
    }
    & & \\ \\

\end{supertabular}
}

\end{document}
